\documentclass[11pt]{article}
\usepackage[margin=1in]{geometry}
\usepackage{booktabs}
\usepackage{longtable}
\usepackage{hyperref}
\usepackage{xcolor}
\usepackage{amsmath}
\usepackage{siunitx}

\title{Replication Report: \emph{Fusobacterium varium} Fv113-g1\\
(Sekizuka et al., 2017, PLOS ONE 12(12):e0189319)}
\author{Ollie (OpenClaw AI) \\ BVBRC Replication Project, Wave BVBRC-100, target \#67}
\date{2026-07-03}

\begin{document}
\maketitle

\begin{abstract}
We independently recompute the descriptive assembly and annotation statistics of the
Sekizuka et al.\ (2017) closed hybrid genome of \emph{Fusobacterium varium} strain Fv113-g1
(RefSeq \texttt{GCF\_002356455.1}, GenBank AP017968--AP017970), plus the comparative
baseline of \emph{F.~varium} ATCC 27725 (\texttt{GCF\_003019655.1}). Eleven descriptive
claims (C1--C11) were checkable from public data; ten reproduce to within annotation-scheme
tolerance and one (FadA paralog count) diverges in a well-understood, pipeline-dependent
way and is marked PARTIAL rather than a contradiction.
\textbf{Verdict: REPLICATED (spot-check level).}
\end{abstract}

\section{Paper}

Sekizuka, Ogasawara, Ohkusa, and Kuroda (2017) report the first closed hybrid assembly of
\emph{F.~varium} Fv113-g1, isolated from a Japanese patient with ulcerative colitis (UC).
Sequencing combined Illumina MiSeq paired-end and mate-pair, PacBio RSII long reads, Argus
optical mapping, and iCORN2 polishing, yielding one \SI{3.96}{\mega\bp} chromosome plus two
plasmids (pFV113-g1-1 and pFV113-g1-2). Annotation with RAST + InterPro + BLASTp identified
an expanded repertoire of putative virulence factors --- \textbf{44 autotransporters (T5SS)}
and \textbf{13 FadA (Fusobacterium adhesin) paralogs} --- positioned as candidate mediators
of mucosal inflammation. Comparative TBLASTx against four other \emph{F.~varium} isolates
(ATCC 8501\textsuperscript{T}, ATCC 27725, ATCC 49185, 12-1B), FastTree phylogeny, and
OrthoVenn orthology argue for a distinct \emph{F.~varium} sub-clade with partial ortholog
sharing with \emph{F.~ulcerans}. RNA-seq under D-MEM vs.\ BHI supports condition-dependent
upregulation of the accessory T5SS and \emph{fadA} paralogs.

\section{Claims tested}

Thirteen headline claims were catalogued; C1--C11 are descriptive/annotation claims
verifiable from the deposited public assembly and were tested in this pass. C12 (RNA-seq
differential expression, DRA005507) and C13 (IS-family enumeration) require re-processing
raw reads or running dedicated repeat finders and are explicitly out of scope for this
spot-check; neither is disputed here.

\section{Method}

\textbf{One layer of evidence: independent recomputation on the paper's own deposited assembly.}

\subsection{Data acquisition}
Three assemblies were pulled with the NCBI Datasets CLI (public, unauthenticated):
\texttt{GCF\_002356455.1} (Fv113-g1), \texttt{GCF\_003019655.1} (\emph{F.~varium} ATCC 27725),
and \texttt{GCF\_037956035.1} (\emph{F.~ulcerans} SB070). Provenance is preserved via each
bundle's \texttt{assembly\_data\_report.jsonl}, \texttt{md5sum.txt}, and NCBI \texttt{README.md}.

\subsection{Recomputation}
Pure Python 3.13 (standard library only), run on 2026-07-03 in the target directory:
(1)~FASTA parse of \texttt{GCF\_002356455.1\_ASM235645v1\_genomic.fna} for per-replicon
length and G+C base counts (N bases treated as non-GC); (2)~GFF feature tally of the RefSeq
PGAP \texttt{genomic.gff} for \texttt{CDS}, \texttt{tRNA}, \texttt{rRNA} counts and pseudogene
flags; (3)~FAA product-name regex scan of \texttt{protein.faa} for
\texttt{FadA}/\texttt{fusobacterium adhesin} and \texttt{autotransporter} matches;
(4)~cross-reference against RefSeq's own \texttt{assemblyStats} and
\texttt{annotationInfo.stats.geneCounts} for internal consistency. No new alignments,
gene predictions, or ortholog searches were performed --- this is an assembly-and-annotation
verification, not a de-novo re-analysis.

\subsection{Tool versions}
Python 3.13 (CPython, Apple silicon); NCBI Datasets CLI (version stamped per bundle);
RefSeq PGAP annotation as most recently run by NCBI on the deposited assembly. The paper's
original annotation used RAST 2.0 + InterPro v49.0 + BLASTp; a $\pm 5\%$ drift on gene
tallies between RAST-2017 and PGAP-current is expected and treated as within-tolerance.

\section{Results vs.\ paper}

\begin{longtable}{clrrrl}
\toprule
\# & Claim & Paper & Measured & $\Delta$ & Verdict \\
\midrule \endhead
C1  & Chromosome length              & \SI{3.96}{\mega\bp}  & \SI{3.965155}{\mega\bp} & $+0.13\%$      & \textbf{MATCH} \\
C2  & Chromosome GC\%                & $29.2\%$             & $29.17\%$               & $-0.03$ pp     & \textbf{MATCH} \\
C3  & Replicons (chr + 2 plasmids)   & $1+2$                & AP017968 + AP017969 + AP017970 & 0     & \textbf{MATCH} \\
C4  & pFV113-g1-1 GC\%               & $26.7\%$             & $26.70\%$               & $0$            & \textbf{MATCH} \\
C5  & pFV113-g1-2 GC\%               & $27.7\%$             & $27.65\%$               & $-0.05$ pp     & \textbf{MATCH} \\
C6  & tRNA gene count                & $58$                 & $58$                    & $0$            & \textbf{MATCH} \\
C7  & rRNA operons                   & $7$                  & $7$ (22 rRNA / 3)       & $0$            & \textbf{MATCH} \\
C8  & Chromosome CDS                 & $3{,}552$            & $3{,}586$ prot-cod / $3{,}671$ CDS & $+0.96\%$ / $+3.4\%$ & \textbf{MATCH (drift)} \\
C9  & ATCC 27725 size                & \SI{3.30}{\mega\bp}  & \SI{3.346458}{\mega\bp} & $+1.4\%$       & \textbf{MATCH (rounding)} \\
C10 & Autotransporter paralogs       & $44$                 & $45$                    & $+1$           & \textbf{MATCH} \\
C11 & FadA paralogs                  & $13$                 & $8$                     & $-5$           & \textcolor{orange}{\textbf{PARTIAL}} \\
\bottomrule
\end{longtable}

\noindent \textbf{Internal consistency.} RefSeq \texttt{assemblyStats} independently reports
\texttt{totalSequenceLength: 4{,}122{,}841}, \texttt{gcPercent: 29},
\texttt{numberOfScaffolds: 3}, and \texttt{geneCounts: \{total: 3754, proteinCoding: 3586,
nonCoding: 83, pseudogene: 85\}}, all agreeing with direct FASTA/GFF measurements and with
the paper.

\section{Notes and caveats}

\begin{itemize}
  \item \textbf{C8 (CDS count)} --- the +34 protein-coding gene delta is a routine
        RAST-2017 vs.\ PGAP-current re-annotation artefact, not a contradiction.
  \item \textbf{C11 (FadA paralog count)} --- a legitimate annotation-scheme divergence.
        The paper's 13 paralogs come from RAST + BLASTp + InterPro homology across the
        broader \emph{fadA}-domain family (likely including hemagglutinin-related
        sub-family paralogs). RefSeq PGAP is more conservative about extending the FadA
        product-name label, yielding only 8 exact matches. An HMMER/Pfam scan against
        PF09403 on \texttt{protein.faa} would arbitrate; the paper's number is not
        falsified by strict product-name counting alone.
  \item \textbf{Abstract-only readers may see a spurious ${\sim}4\%$ size discrepancy}
        because the abstract says ``3.96 Mb'' (chromosome only) while the RefSeq
        \texttt{totalSequenceLength} 4.12 Mb includes the two plasmids
        ($\SI{89.6}{\kilo\bp} + \SI{68.1}{\kilo\bp}$). Table 1 of the paper is consistent
        with the whole-genome total.
  \item \textbf{C12 / C13} --- RNA-seq DE (DRA005507) and IS-family enumeration are
        skipped in this pass. Both are testable from public data but were deliberately
        deferred; neither is disputed.
  \item \textbf{Comparator drift (C9).} The 1.4\% size delta between the paper's
        ATCC 27725 draft (2017) and the current complete RefSeq assembly is expected and
        does not change the qualitative comparative claim that Fv113-g1 is ${\sim}600$ kb
        larger than ATCC 27725.
\end{itemize}

\section{Genuine critique}

This section deliberately reports the honest weaknesses of \emph{this replication effort}
(not of the paper). It is included so that downstream reviewers can distinguish
``we checked everything and it held'' from ``we checked the tractable subset''.

\begin{enumerate}
  \item \textbf{Single evidence layer.} The verification is one layer deep: direct
        recomputation on the deposited assembly. We did not independently redo the hybrid
        assembly from raw Illumina + PacBio reads, did not rerun RAST/InterPro, and did
        not rerun OrthoVenn/FastTree. A conclusion of ``REPLICATED'' therefore covers
        \emph{the derivative statistics on the deposited artefact}, not the assembly
        pipeline itself.
  \item \textbf{Central biological claims C12 and C13 are untested.} The differential
        expression signal for T5SS/\emph{fadA} paralogs in D-MEM vs.\ BHI (C12) and the
        ISFv1/ISFv2 enumerations (C13) are the paper's more mechanism-relevant claims and
        were deferred. This is a limitation of scope, not a critique of the paper.
  \item \textbf{FadA count divergence unresolved.} C11 is marked PARTIAL, but the
        arbitrating experiment (HMMER against PF09403, or a curated \emph{fadA} MSA-based
        profile search) was not performed. We report the divergence honestly and do not
        claim to have vindicated either number.
  \item \textbf{Product-name string matching is fragile.} Both C10 (+1) and C11 ($-5$)
        rest on case-insensitive regex matches against PGAP product names. A single
        product-string wording change in a future PGAP re-annotation would silently
        shift both numbers. A domain-model-based paralog count would be more durable.
  \item \textbf{Comparator scope is minimal.} Only one comparator strain (ATCC 27725) was
        checked for C9. The paper compares against four other \emph{F.~varium} strains
        plus multiple \emph{F.~ulcerans} genomes; those additional comparisons were not
        re-verified.
  \item \textbf{PGAP version drift is not pinned.} We record ``RefSeq PGAP as most
        recently run'' but do not pin the exact PGAP release. A stricter replication would
        capture the PGAP version and re-run identically against a frozen release.
  \item \textbf{No orthogonal comparator on GC/length.} The internal consistency check is
        derived from the same underlying FASTA (assemblyStats vs.\ our own parse). Two
        digestions of the same file are not a truly independent second layer.
  \item \textbf{Ambient bias risk.} An observer with prior knowledge that a paper's
        deposited assembly should ``match itself'' is unlikely to be surprised by a
        matching result. We caution readers that a positive spot-check on deposited
        artefacts is a \emph{necessary but not sufficient} condition for full replication.
\end{enumerate}

\section{Verdict}

\textbf{REPLICATED (spot-check).} Every checkable descriptive claim (C1--C10) matches to
within measurement or annotation-drift tolerance. C11 is PARTIAL (annotation-scheme
divergence), not CONTRADICTED. The paper's higher-level interpretive claims are not
re-tested here but rest on an assembly that is verifiably faithful to the deposited public
record.

\section{Files}
\begin{itemize}
  \item \texttt{report/REPORT.md}, \texttt{report/REPORT.tex}
  \item \texttt{report/evidence/fv113g1\_assembly\_stats.json}
  \item \texttt{report/evidence/comparative\_genomes.json}
  \item \texttt{report/evidence/claims\_vs\_measured.csv}
  \item \texttt{work/data/\{Fv113g1,ATCC27725,Fulcerans\}/} --- staged NCBI Datasets bundles
  \item \texttt{work/esearch\_*.json}, \texttt{work/esummary\_asm.json} --- accession provenance
\end{itemize}

\vspace{1em}
\noindent\emph{Sekizuka et al.\ 2017 · PLOS ONE 12(12): e0189319 · GCF\_002356455.1 ·
Verdict: REPLICATED · Fv113-g1 chromosome \num{3965155} bp / 29.17\% GC / 58 tRNA /
7 rRNA operons --- matches paper.}

\end{document}
